\section{Executable Runtime Model}
\label{sec:model}

The model is intentionally small, but it is not a claim of expressive
minimality.  Its purpose is to make replay, scheduling, and effects observable
enough to falsify claimed properties.

\subsection{Events, projection, and identity}

An event envelope is a tuple

\begin{equation}
 e = (id, run, seq, kind, cause, emitter).
 \label{eq:event-envelope}
\end{equation}

The event kind is a closed sum containing domain facts, signals, effect
lifecycle events, and evidence facts.  The evidence branch includes
\code{HazardDetected(detail)}; an envelope records it as
\code{EvidenceRecorded(HazardDetected detail)}, and the grader treats that fact
as an explicit blocker.  A log
$L = \langle e_1,\ldots,e_n\rangle$ is ordered and append-only.  The pure
reducer $\mathsf{evolve}:S\times E\rightarrow S$ yields

\begin{equation}
 \project(L) = \mathsf{fold}(\mathsf{evolve}, S_0, L).
 \label{eq:project}
\end{equation}

Event identity is idempotent at the reducer boundary: if an event identifier
has already been applied, a repeated envelope does not mutate state again.  A
duplicate identifier remains an integrity fault for fork assessment, because
idempotent reduction must not silently certify a malformed trace.

State contains domain facts and signals, projected effect descriptors,
evidence facts, integrity faults, and the set of applied event identifiers.
When comparing projected domain state, the implementation excludes the
bookkeeping set of applied identifiers.  This prevents a replay that merely
renames envelopes from being declared semantically different while still
allowing exact-trace equality to detect the change.

\begin{proposition}[Closed-log projection determinism]
For a fixed ordered event log and a pure reducer, repeated evaluation of
Equation~\ref{eq:project} returns the same state.  Oracle calls are not an
exception: their \emph{recorded results}, when present, are events in the fixed
log.
\end{proposition}

This proposition is the trivial base case.  It says nothing about how the log
was produced, whether the event order was unique, whether missing oracle
results can be recovered, or whether effects may be executed again.

\subsection{Behaviors and actions}

A behavior is a tuple

\begin{equation}
 b=(id_b, trigger_b, fire_b, W_b)
\end{equation}

where $trigger_b:S\times E\rightarrow\mathbb{B}$ decides whether an event
enables the behavior, $fire_b:S\times E\rightarrow A^*$ returns actions, and
$W_b$ is a declared write set used only as a diagnostic.  Actions can emit a
domain fact or signal, or request an external effect.  Behaviors cannot emit an
effect outcome or evidence attestation directly; those facts must cross the
interpreter or validation boundary.  This prevents a reducer from minting its
own proof that an external action succeeded or that a run was verified.

The kernel does not execute actions.  Emissions become pending events.  Effect
requests become pending request events and later block until an interpreter
injects a committed, failed, or unknown outcome.  Ordinary .NET tasks, network
clients, mutable caches, and database adapters may exist outside the boundary;
none are reachable from $trigger$ or $fire$.

\subsection{Configurations and scheduling}

A configuration contains projected state, pending events, enabled activations,
outstanding effects, the retained trace, the behavior firing trace, and the next
sequence and generated-identity counters.  A scheduler has two explicit
selectors: one chooses a pending event, and the other chooses an enabled
activation.  Four policies are included: canonical, reverse-event,
reverse-activation, and both reversed.

Writing a configuration as $\langle s,Q,A,O,L,T\rangle$, the two load-bearing
transition rules are shown in Figure~\ref{fig:small-step}.  The helper
$\mathsf{materialize}$ turns a behavior's pure action descriptions into pending
fact/signal events or effect-request events.  It never produces an interpreter
outcome.

\begin{figure}[ht]
\centering
\small
\begin{minipage}{0.96\textwidth}
\[
\frac{
 e=\mathsf{selectEvent}_{\sigma}(Q)\quad
 s'=\mathsf{evolve}(s,e)\quad
 A'=A\mathbin{+}\{(b,e)\mid \mathsf{trigger}_b(s',e)\}
}{
 \langle s,Q,A,O,L,T\rangle
 \rightarrow_{\sigma}
 \langle s',Q\setminus e,A',O,L\mathbin{+}e,T\rangle
}
\quad\textsc{Process-Event}
\]
\[
\frac{
 a=(b,e)=\mathsf{selectActivation}_{\sigma}(A)\quad
 X=\mathsf{fire}_b(s,e)\quad
 (Q',O')=\mathsf{materialize}(Q,O,X)
}{
 \langle s,Q,A,O,L,T\rangle
 \rightarrow_{\sigma}
 \langle s,Q',A\setminus a,O',L,T\mathbin{+}(b,e)\rangle
}
\quad\textsc{Fire-Activation}
\]
\end{minipage}
\caption{Small-step scheduling rules.  Event and activation selection are
explicit.  Emissions become visible to state only after a later
\textsc{Process-Event} step, not within the activation batch for one trigger.}
\label{fig:small-step}
\end{figure}

Processing a pending event first applies it to state and then computes enabled
behaviors from the updated state.  Enabled behaviors fire sequentially, but
their emissions enter the pending queue.  Activations handling the same
triggering event therefore observe the same projected state.  A behavior
activated by a later-processed event observes changes committed by earlier
\textsc{Process-Event} steps, including events materialized from prior
activations.  This cross-event ordering seam mirrors the one found in
ActiveGraph and is sufficient for the read/trigger interference witness in
Section~\ref{sec:confluence}.

Let $C\rightarrow_\sigma C'$ denote one legal event-processing or
activation-firing step under scheduler $\sigma$.  Settlement repeatedly steps
until one of three outcomes is reached:

\begin{itemize}[leftmargin=*]
  \item \textbf{Quiescent}: no pending events, enabled activations, or
  outstanding effects;
  \item \textbf{BlockedAwaitingEffect}: no internal work remains, but at least
  one request has no terminal outcome;
  \item \textbf{Diverged}: a caller-supplied step bound is reached before either
  terminal outcome.
\end{itemize}

The step bound is checked only after testing for a terminal configuration.  A
run that becomes quiescent on the exact final step is therefore not
misclassified as divergence.  Conversely, reaching the bound does not prove a
cycle; the kernel reports \code{StepBoundReached} and makes no stronger claim.
BlockedAwaitingEffect is terminal for the settlement procedure but is not a
quiescent normal form.  The confluence definition in Section~\ref{sec:confluence}
ranges only over executions that reach Quiescent; uniqueness under blocking is
a different property and is not claimed here.

\subsection{Observations and equivalence}

Confluence and fork claims are meaningless without specifying what is observed.
The artifact implements three equivalences:

\begin{enumerate}[leftmargin=*]
  \item \textbf{ExactTrace}: envelope and event lists are structurally equal;
  \item \textbf{ProjectedState}: domain projection is equal after excluding
  applied-event bookkeeping and envelope metadata;
  \item \textbf{FactsOnly}$(K)$: the selected fact keys agree.
\end{enumerate}

The validation additionally reports byte equality when a source artifact
provides it, decision equality when the study records decisions, and structural
prefix equality modulo explicitly excluded lineage fields.  These relations
must not be substituted for one another.  In particular, equal decisions do
not imply equal paths, and equal projected state does not imply an unchanged
external world.
